% A3 试卷示例：exam-zh，A3 横版双栏 + 密封线 + 卷头
\documentclass{exam-zh}
\examsetup{
  font = termes, math-font = termes,
  page = { size = a3paper },
  sealline = { show = true, scope = everypage },
}
\title{2026 年秋季高一数学期中考试}
\subject{数学}
\begin{document}
\maketitle
\begin{notice}
\item 本卷共 2 页，15 题．全卷满分 120 分．考试用时 100 分钟．
\item 答题前，先将自己的姓名、准考证号、考场号、座位号填写在试卷和答题卡上．
\end{notice}
\information{姓名\underline{\hspace{6em}}，座位号\underline{\hspace{6em}}}

\section{选择题（每题 5 分，共 40 分）}
\begin{question}
已知集合 $A=\{1,2,3\}$，$B=\{2,3,4\}$，则 $A\cap B=$ \paren
\begin{choices}
\item $\{1,2\}$
\item $\{2,3\}$
\item $\{3,4\}$
\item $\{1,4\}$
\end{choices}
\end{question}
\begin{question}
函数 $f(x)=x^2-2x$ 的最小值为 \paren
\begin{choices}
\item $-2$
\item $-1$
\item $0$
\item $1$
\end{choices}
\end{question}
\begin{question}
若 $a>b>0$，则下列不等式一定成立的是 \paren
\begin{choices}
\item $a^2<b^2$
\item $\dfrac{1}{a}>\dfrac{1}{b}$
\item $ac>bc$
\item $a+c>b+c$
\end{choices}
\end{question}
\begin{question}
已知 $\sin\alpha=\dfrac{3}{5}$，$\alpha\in\left(0,\dfrac{\pi}{2}\right)$，则 $\cos\alpha=$ \paren
\begin{choices}
\item $\dfrac{4}{5}$
\item $\dfrac{3}{4}$
\item $\dfrac{5}{4}$
\item $\dfrac{2}{5}$
\end{choices}
\end{question}
\begin{question}
复数 $z=1+i$，则 $z^2=$ \paren
\begin{choices}
\item $2i$
\item $-2i$
\item $2$
\item $1-i$
\end{choices}
\end{question}
\begin{question}
已知向量 $\vec a=(1,2)$，$\vec b=(3,-1)$，则 $\vec a\cdot\vec b=$ \paren
\begin{choices}
\item $1$
\item $-1$
\item $5$
\item $7$
\end{choices}
\end{question}
\begin{question}
函数 $y=\log_2(x-1)$ 的定义域为 \paren
\begin{choices}
\item $(0,+\infty)$
\item $(1,+\infty)$
\item $[1,+\infty)$
\item $\mathbf{R}$
\end{choices}
\end{question}
\begin{question}
已知等差数列 $\{a_n\}$ 中 $a_1=1$，公差 $d=2$，则 $a_{10}=$ \paren
\begin{choices}
\item $19$
\item $20$
\item $21$
\item $22$
\end{choices}
\end{question}

\section{填空题（每题 5 分，共 20 分）}
\begin{question}
计算 $\displaystyle\int_0^1 x^2\,\mathrm{d}x=$ \fillin[$\dfrac{1}{3}$]
\end{question}
\begin{question}
已知 $\tan\alpha=2$，则 $\dfrac{\sin\alpha+\cos\alpha}{\sin\alpha-\cos\alpha}=$ \fillin[$3$]
\end{question}
\begin{question}
曲线 $y=x^3-3x$ 在点 $(1,-2)$ 处的切线方程为 \fillin[$y=-2$]
\end{question}
\begin{question}
$(1+x)^6$ 的展开式中 $x^3$ 的系数为 \fillin[$20$]
\end{question}

\section{解答题（共 60 分）}
\begin{problem}
（20 分）已知函数 $f(x)=x^3-3x^2+2$。
\begin{enumerate}
\item 求 $f(x)$ 的单调区间；
\item 求 $f(x)$ 在 $[-1,3]$ 上的最大值与最小值。
\end{enumerate}
\end{problem}
\begin{solution}
$f'(x)=3x^2-6x=3x(x-2)$。当 $x<0$ 或 $x>2$ 时 $f'(x)>0$；当 $0<x<2$ 时 $f'(x)<0$。
故递增区间为 $(-\infty,0)$ 与 $(2,+\infty)$，递减区间为 $(0,2)$。
又 $f(-1)=-2$，$f(0)=2$，$f(2)=-2$，$f(3)=2$，所以最大值为 $2$，最小值为 $-2$。
\end{solution}
\begin{problem}
（20 分）在 $\triangle ABC$ 中，角 $A,B,C$ 的对边分别为 $a,b,c$，已知 $a=2$，$b=3$，$C=\dfrac{\pi}{3}$。
\begin{enumerate}
\item 求 $c$；
\item 求 $\triangle ABC$ 的面积。
\end{enumerate}
\end{problem}
\begin{solution}
由余弦定理 $c^2=a^2+b^2-2ab\cos C=4+9-6=7$，故 $c=\sqrt{7}$。
面积 $S=\dfrac12 ab\sin C=\dfrac12\cdot2\cdot3\cdot\dfrac{\sqrt3}{2}=\dfrac{3\sqrt3}{2}$。
\end{solution}
\begin{problem}
（20 分）已知椭圆 $\dfrac{x^2}{a^2}+\dfrac{y^2}{b^2}=1$（$a>b>0$）的离心率为 $\dfrac{\sqrt3}{2}$，且过点 $\left(1,\dfrac{\sqrt3}{2}\right)$，求椭圆方程。
\end{problem}

\end{document}
